% =========================================================
% Section: Chain Rule
% =========================================================
% ---------------------------------------------------------

\begin{theorembox}{Theorem}
\begin{theorem}[Chain Rule]
\label{thm:chain-rule}
Let $I, J \subseteq \mathbb{R}$ be intervals, let $f : I \to \mathbb{R}$ and $g : J \to \mathbb{R}$, and suppose that $f(I) \subseteq J$. Let $c \in I$. If $f$ is differentiable at $c$ and $g$ is differentiable at $f(c)$, then the composition $g \circ f$ is differentiable at $c$, and
\[
(g \circ f)'(c) = g'(f(c))\,f'(c).
\]
\hyperref[prf:chain-rule]{Go to proof.}
\end{theorem}
\end{theorembox}

\begin{remark*}[Standard quantified statement]
\[
  \operatorname{IsDifferentiable}(f,c)
  \land
  \operatorname{IsDifferentiable}(g,f(c))
  \Longrightarrow
  (g\circ f)'(c)=g'(f(c))f'(c).
\]
\end{remark*}

\begin{remark*}[Interpretation]
The derivative of a composition is the derivative of the outside function evaluated at the inside value, multiplied by the derivative of the inside function.


Let $\phi_f$ and $\phi_g$ be the Carath\'{e}odory slope factors for $f$ at $c$
and $g$ at $f(c)$. Then
\[
g(f(x))-g(f(c))
  =
\phi_g(f(x))(f(x)-f(c))
  =
\phi_g(f(x))\phi_f(x)(x-c).
\]
The new slope factor is
\[
  \phi_{g\circ f}(x)=\phi_g(f(x))\phi_f(x),
\]
which is continuous at $c$ and has value $g'(f(c))f'(c)$. This proof never
divides by $f(x)-f(c)$, so it avoids the failure in the informal
$\Delta g/\Delta f \cdot \Delta f/\Delta x$ argument when $f(x)=f(c)$ near
$c$.
\end{remark*}

\begin{dependencies}
\begin{itemize}
  \item \hyperref[thm:caratheodory-characterization-of-differentiability]{Theorem: Carath\'{e}odory Characterization of Differentiability}
  \item \hyperref[thm:differentiable-implies-continuous]{Theorem: Differentiable Implies Continuous}
  \item \hyperref[thm:composition-of-limits]{Theorem: Composition of Limits}
  \item \hyperref[thm:limit-product]{Theorem: Product Rule for Function Limits}
\end{itemize}
\end{dependencies}

\begin{theorembox}{Theorem (Leibniz Rule)}
\begin{theorem}[Leibniz Rule]
\label{thm:leibniz-rule}
Let \(I\subseteq\mathbb{R}\) be an interval, let \(x\in I\), and suppose
that \(f,g:I\to\mathbb{R}\) are \(n\)-times differentiable on a
neighbourhood of \(x\). Then \(fg\) is \(n\)-times differentiable at \(x\),
and
\[
  (fg)^{(n)}(x)
  =
  \sum_{k=0}^{n}\binom{n}{k}f^{(k)}(x)g^{(n-k)}(x).
\]
\hyperref[prf:leibniz-rule]{Go to proof.}
\end{theorem}
\end{theorembox}

\begin{remark*}[Standard quantified statement]
\[
  f,g\text{ are }n\text{-times differentiable at }x
  \Longrightarrow
  (fg)^{(n)}(x)
  =
  \sum_{k=0}^{n}\binom{n}{k}f^{(k)}(x)g^{(n-k)}(x).
\]
\end{remark*}

\begin{remark*}[Interpretation]
The binomial pattern is not an accident. Each of the \(n\) differentiations
lands on either \(f\) or \(g\), and \(\binom{n}{k}\) counts the ways to send
exactly \(k\) differentiations to \(f\). At \(n=1\), this is the ordinary
product rule.
\end{remark*}

\begin{dependencies}
\begin{itemize}
  \item \hyperref[thm:product-rule]{Theorem: Product Rule}
  \item \hyperref[def:higher-derivatives]{Definition: Higher Derivatives}
\end{itemize}
\end{dependencies}

\begin{theorembox}{Theorem (Faa di Bruno's Formula)}
\begin{theorem}[Faa di Bruno's Formula]
\label{thm:faa-di-bruno}
Let \(I,J\subseteq\mathbb{R}\) be intervals, let \(x\in I\), suppose
\(g:I\to J\) is \(n\)-times differentiable on a neighbourhood of \(x\), and
suppose \(f:J\to\mathbb{R}\) is \(n\)-times differentiable on a neighbourhood
of \(g(x)\). Then \(f\circ g\) is \(n\)-times differentiable at \(x\), and
\[
  \frac{d^n}{dx^n}f(g(x))
  =
  \sum_{\sum_{m=1}^{n}mk_m=n}
  \frac{n!}{\prod_{m=1}^{n}k_m!\,(m!)^{k_m}}\,
  f^{(k)}(g(x))
  \prod_{m=1}^{n}\bigl(g^{(m)}(x)\bigr)^{k_m},
\]
where \(k_m\geq0\) and \(k=\sum_{m=1}^{n}k_m\).
\hyperref[prf:faa-di-bruno]{Go to proof.}
\end{theorem}
\end{theorembox}

\begin{remark*}[Standard quantified statement]
\[
  g\text{ is }n\text{-times differentiable at }x
  \land
  f\text{ is }n\text{-times differentiable at }g(x)
  \Longrightarrow
  (f\circ g)^{(n)}(x)
  =
  \sum_{\sum mk_m=n}
  \frac{n!}{\prod_m k_m!\,(m!)^{k_m}}\,
  f^{(k)}(g(x))\prod_m\bigl(g^{(m)}(x)\bigr)^{k_m}.
\]
\end{remark*}

\begin{remark*}[Interpretation]
The chain rule gives the first derivative of a composition. Repeated
differentiation is deeper: product-rule terms and chain-rule terms interlace,
and the bookkeeping is governed by integer partitions. Faa di Bruno's
formula is the all-orders chain rule.


Smoothness is more than this formula needs. The \(n\)-th order statement
requires \(n\) derivatives of \(g\) near \(x\) and \(n\) derivatives of \(f\)
near \(g(x)\). At \(n=1\), the formula is exactly the chain rule.


The tuple \((k_1,\ldots,k_n)\) records an integer partition of \(n\): \(k_m\)
is the number of parts of size \(m\). The coefficient
\[
  \frac{n!}{\prod_m k_m!\,(m!)^{k_m}}
\]
counts set partitions with that block-size profile. The outer derivative
\(f^{(k)}\) is differentiated once per block, while each block of size \(m\)
contributes one factor \(g^{(m)}\).


The formula uses the bare-product normalization: the coefficient contains
\((m!)^{k_m}\), and the product contains \(\bigl(g^{(m)}(x)\bigr)^{k_m}\).
Equivalently, one may put \(g^{(m)}(x)/m!\) inside the product and remove
\((m!)^{k_m}\) from the denominator. Mixing these two conventions gives the
wrong coefficients.


The constraint \(\sum_m mk_m=n\) is precisely the multiplicity form of an
integer partition of \(n\). Thus the number of summands is \(p(n)\), the
partition function from the discrete and number-system material, now
appearing as bookkeeping for iterated differentiation.


This theorem is enrichment, not a critical dependency for the integration,
measure, or manifold arcs. It is included because it is the complete
all-orders chain rule and a useful callback to partitions.
\end{remark*}

\begin{dependencies}
\begin{itemize}
  \item \hyperref[thm:chain-rule]{Theorem: Chain Rule}
  \item \hyperref[thm:leibniz-rule]{Theorem: Leibniz Rule}
  \item \hyperref[def:higher-derivatives]{Definition: Higher Derivatives}
  \item \hyperref[def:class-ck]{Definition: Class \(C^k\)}
\end{itemize}
\end{dependencies}

\begin{exposition}
\phantomsection\label{ex:faa-di-bruno-n2}
The partitions of \(2\) are represented by \((k_1,k_2)=(2,0)\) and
\((0,1)\). The first gives \(f''(g)(g')^2\), and the second gives
\(f'(g)g''\). Hence
\[
  (f\circ g)''
  =
  f''(g)(g')^2+f'(g)g''.
\]
The coefficient on \(f'(g)g''\) is \(1\), which is the quickest check that
the normalization has not been mixed.
\end{exposition}

\begin{exposition}
\phantomsection\label{ex:faa-di-bruno-n3}
The partitions of \(3\) produce the three terms
\[
  (f\circ g)'''
  =
  f'''(g)(g')^3
  +3f''(g)g'g''
  +f'(g)g'''.
\]
The middle coefficient \(3\) counts the three ways to choose which
differentiation contributes the second-derivative block of \(g\).
\end{exposition}
